\documentclass[11pt,letterpaper]{article}

\usepackage[margin=0.65in, top=0.55in, bottom=0.55in]{geometry}
\usepackage{iftex}
\ifPDFTeX
    \usepackage[utf8]{inputenc}
    \usepackage[T1]{fontenc}
\fi
\usepackage{hyperref}
\usepackage[usenames,dvipsnames]{xcolor}
\usepackage{enumitem}
\usepackage{titlesec}
\usepackage{fancyhdr}
\usepackage{multicol}
\usepackage{fontawesome}
\usepackage{academicons} % real Google Scholar glyph (needs lualatex/xelatex)

% Google Scholar icon: true academicons logo on lua/xelatex, graceful fallback on pdflatex
\ifPDFTeX
    \newcommand{\scholaricon}{\faGraduationCap}
\else
    \newcommand{\scholaricon}{\aiGoogleScholar}
\fi

% Colors
\definecolor{linkblue}{HTML}{0B5394}

\hypersetup{
    colorlinks=true,
    linkcolor=linkblue,
    urlcolor=linkblue
}

% Section formatting
\titleformat{\section}{\large\bfseries\scshape}{}{0em}{}[\titlerule]
\titlespacing{\section}{0pt}{8pt}{4pt}

% Remove page numbering for single page or add custom
\pagestyle{fancy}
\fancyhf{}
\renewcommand{\headrulewidth}{0pt}
\fancyfoot[C]{\footnotesize Last updated: July 2026}

% Tight lists
\setlist[itemize]{nosep, leftmargin=24pt, topsep=2pt}

% Compact dated entries for teaching and talks
\newcommand{\cvdatedentry}[3]{%
    \noindent
    \parbox[t]{0.75\textwidth}{%
        \raggedright\textbf{#1}\\[-1pt]
        {\small #2}\par
    }%
    \hfill
    \parbox[t]{0.22\textwidth}{%
        \raggedleft{\small #3}\par
    }%
    \par\vspace{2pt}
}

\begin{document}

% === HEADER ===
\begin{center}
    {\LARGE \textbf{Dong-Kyum Kim}} \\[6pt]
    \href{mailto:dong-kyum.kim@mpi-sp.org}{\faEnvelope\,\, dong-kyum.kim@mpi-sp.org} $\cdot$
    \href{https://kdkyum.github.io}{\faGlobe\,\, kdkyum.github.io} $\cdot$
    \href{https://x.com/kdkyum531}{\faTwitter\,\, @kdkyum531} $\cdot$
    \href{https://scholar.google.com/citations?user=-pvD9xUAAAAJ}{\scholaricon\,\, Google Scholar} $\cdot$
    \href{https://www.linkedin.com/in/kdkyum}{\faLinkedin\,\, LinkedIn} \\[5pt]
    {\small \faMapMarker\,\, Max Planck Institute for Security and Privacy, Universit\"atsstra\ss e 140, 44799 Bochum, Germany}
\end{center}

% === RESEARCH INTERESTS ===
\section{Research Interests}
\begin{itemize}
    \item \textbf{Interdisciplinary physicist and AI researcher}, guided by physical principles and observations from neuroscience.
    \item \textbf{Knowledge in large language models}: uncovering how models represent, retrieve, update, and forget information through mechanistic interpretability, model editing, and machine unlearning.
    \item \textbf{From physical and biological systems to AI}: drawing on nonequilibrium statistical physics and biological learning and memory to build safer, more reliable, and controllable AI systems.
\end{itemize}

% === EDUCATION ===
\section{Education}
\noindent\textbf{Korea Advanced Institute of Science and Technology (KAIST)} \hfill 2016--2022 \\
Ph.D. in Physics \hfill Advisor: Prof.\ \href{https://stat.kaist.ac.kr:2583/~hjeong/}{Hawoong Jeong} \\
{\small Dissertation: \textit{Nonequilibrium Statistical Physics Study using Deep Learning}}

\vspace{4pt}
\noindent\textbf{Seoul National University (SNU)} \hfill 2011--2015 \\
B.S. in Physics, Minor in Computer Science \& Engineering

% === EXPERIENCE ===
\section{Research Experience}

\noindent\textbf{Postdoctoral Researcher}, Max Planck Institute for Security and Privacy (MPI-SP) \hfill Oct.\
2024--Present
\begin{itemize}
    \item Hosted by Prof.\ \href{https://www.linkedin.com/in/miacha/}{Meeyoung Cha} (Scientific Director)
    \item Data Science for Humanity Group
\end{itemize}

\vspace{4pt}
\noindent\textbf{Postdoctoral Researcher}, Institute for Basic Science (IBS) \hfill Mar.\ 2022--Sep.\ 2024
\begin{itemize}
    \item Hosted by Prof.\ \href{https://www.linkedin.com/in/miacha/}{Meeyoung Cha} (Chief Investigator)
    \item Data Science Group, Center for Mathematical and Computational Science
\end{itemize}

\vspace{4pt}
\noindent\textbf{Data Science Intern}, Samsung Electronics \hfill Sep.--Dec.\ 2017
\begin{itemize}
    \item Built distributed ML systems with \href{https://www.linkedin.com/in/danielyounghokim/}{Daniel Kim} (Senior Data Scientist): multi-GPU anomaly image classification (Keras \& Spark) and an Elasticsearch-based image search framework
\end{itemize}

% === PUBLICATIONS ===
\section{Publications \hfill {\normalfont\small $^\dagger$ denotes equal contribution.}}

\begin{enumerate}[label={[\arabic*]}, leftmargin=20pt, itemsep=3pt]

    \item Jea Kwon, \textbf{Dong-Kyum Kim}, Jiwon Kim, Yonghyun Kim, Woong Kook, and Meeyoung Cha.
    ``\href{https://arxiv.org/abs/2606.14997}{AI Engram: In Search of Memory Traces in Artificial Intelligence}.''
    \textit{ICML} 2026 \textbf{(Oral)}.

    \item Minsung Kim, \textbf{Dong-Kyum Kim}, Jea Kwon, Nakyeong Yang, Kyomin Jung, and Meeyoung Cha.
    ``\href{https://arxiv.org/abs/2510.02370}{How Training Data Shapes the Use of Parametric and In-Context Knowledge in
Language Models}.''
    \textit{ACL} 2026.

    \item \textbf{Dong-Kyum Kim}, Minsung Kim, Jea Kwon, Nakyeong Yang, and Meeyoung Cha.
    ``\href{https://arxiv.org/abs/2509.21993}{Bilinear representation mitigates reversal curse and enables consistent
model editing}.''
    \textit{ICLR} 2026.

    \item Nakyeong Yang, \textbf{Dong-Kyum Kim}, Jea Kwon, Minsung Kim, Kyomin Jung, and Meeyoung Cha.
    ``\href{https://arxiv.org/abs/2509.22263}{Erase or Hide? Suppressing Spurious Unlearning Neurons for Robust
Unlearning}.''
    \textit{ICLR} 2026.

    \item Jea Kwon, Jiwon Kim, \textbf{Dong-Kyum Kim}, and Meeyoung Cha.
    ``\href{https://openreview.net/forum?id=SHqpI6E5cw}{Moir: Let the Model Direct Its Own Story for Robust Cross-Domain
Knowledge Editing}.''
    \textit{ICML 2026 Workshop on Mechanistic Interpretability}.

    \item Minsung Kim, Jea Kwon, \textbf{Dong-Kyum Kim}, and Meeyoung Cha.
    ``\href{https://d2jud02ci9yv69.cloudfront.net/2025-04-28-engram-172/blog/engram/}{In Search of the Engram in
LLMs}.''
    \textit{ICLR 2025 Blogpost}.

    \item Gwangsu Kim, \textbf{Dong-Kyum Kim}, and Hawoong Jeong.
    ``\href{https://doi.org/10.1038/s41467-023-44516-0}{Spontaneous emergence of rudimentary music detectors in deep
neural networks}.''
    \textit{Nature Communications} 2024.

    \item Jea Kwon, Sunpil Kim, \textbf{Dong-Kyum Kim}, Jinhyeong Joo, SoHyung Kim, Meeyoung Cha, and C. Justin Lee.
    ``\href{https://doi.org/10.1007/s11263-024-02072-0}{SUBTLE: An unsupervised platform with temporal link embedding
that maps animal behavior}.''
    \textit{IJCV} (2024).

    \item \textbf{Dong-Kyum Kim}$^\dagger$, Jea Kwon$^\dagger$, Meeyoung Cha, and C. Justin Lee.
    ``\href{https://openreview.net/forum?id=vKpVJxplmB}{Transformer as a hippocampal memory consolidation model based on
NMDAR-inspired nonlinearity}.''
    \textit{NeurIPS} 2023.

    \item Sangyun Lee, \textbf{Dong-Kyum Kim}, Jong-Min Park, Won Kyu Kim, Hyunggyu Park, and Jae Sung Lee.
    ``\href{https://doi.org/10.1103/PhysRevResearch.5.013194}{Multidimensional entropic bound}.''
    \textit{Phys.\ Rev.\ Research} \textbf{5}, 013194 (2023).

    \item \textbf{Dong-Kyum Kim}$^\dagger$, Sangyun Lee$^\dagger$, and Hawoong Jeong.
    ``\href{https://doi.org/10.1103/PhysRevResearch.4.023051}{Estimating entropy production with odd-parity state
variables via ML}.''
    \textit{Phys.\ Rev.\ Research} \textbf{4}, 023051 (2022).

    \item Youngkyoung Bae, \textbf{Dong-Kyum Kim}, and Hawoong Jeong.
    ``\href{https://doi.org/10.1103/PhysRevResearch.4.033094}{Inferring dissipation maps from videos using CNNs}.''
    \textit{Phys.\ Rev.\ Research} \textbf{4}, 033094 (2022).

    \item \textbf{Dong-Kyum Kim} and Hawoong Jeong.
    ``\href{https://doi.org/10.1103/PhysRevResearch.3.L022002}{Deep reinforcement learning for feedback control in a
collective flashing ratchet}.''
    \textit{Phys.\ Rev.\ Research} \textbf{3}, L022002 (2021).

    \item \textbf{Dong-Kyum Kim}$^\dagger$, Youngkyoung Bae$^\dagger$, Sangyun Lee, and Hawoong Jeong.
    ``\href{https://doi.org/10.1103/PhysRevLett.125.140604}{Learning Entropy Production via Neural Networks}.''
    \textit{Phys.\ Rev.\ Lett.} \textbf{125}, 140604 (2020).

    \item \textbf{Dong-Kyum Kim}$^\dagger$, Byunghwee Lee$^\dagger$, Daniel Kim, and Hawoong Jeong.
    ``\href{https://doi.org/10.3938/jkps.76.368}{Multi-label classification of historical documents using hierarchical
attention}.''
    \textit{J.\ Korean Phys.\ Soc.} \textbf{76}, 368 (2020).

\end{enumerate}

% === TEACHING ===
\section{Teaching Experience}

\cvdatedentry{Computational Physics Course, KAIST (Daejeon, Korea)}
    {\textit{``Deep learning applications: Nonequilibrium statistical physics study using AI''}}
    {May 1, 2023}

\cvdatedentry{Computational Physics Course, KAIST (Daejeon, Korea)}
    {\textit{``Deep Learning Introduction''}}
    {Apr. 24, 2023}

\cvdatedentry{SNU Physics and AI Winter School (Seoul, Korea)}
    {\textit{``Exploring Irreversibility via Machine Learning''}}
    {Feb. 24, 2022}

\cvdatedentry{General Physics II, KAIST (Korea)}
    {Teaching Assistant}
    {Fall 2016; Spring 2017}

% === TALKS ===
\section{Talks}

\cvdatedentry{Spark(l)ing Science, MPI-SP (MB/1 Seminar Room + Zoom)}
    {\href{https://docs.google.com/presentation/d/1WdZvexF7HmzqnJbV-Xiu-dhJM0EwJcDLm17H_DoYrzE/edit?usp=sharing}{\textit{``AI research using agents and notebooks''}}}
    {Jul. 24, 2026}

\cvdatedentry{\href{https://sites.google.com/view/cxrex/home}{CREx \#23} (Online)}
    {\textit{``Bilinear relational structure fixes reversal curse and enables consistent model editing''}}
    {Feb. 24, 2026}

\cvdatedentry{IBS Winter School on AI-Boosted Basic Science (Daejeon, Korea)}
    {\textit{``Resemblances between Transformer's Nonlinearity and NMDA Receptor Dynamics''}}
    {Dec. 13, 2022}

\cvdatedentry{KIAS CAINS Summer Workshop (Jeju, Korea)}
    {\textit{``Working and reference memory in transformers on a navigation task''}}
    {Sep. 2, 2022}

\cvdatedentry{KIAS Nonequilibrium Statistical Physics of Complex Systems (Seoul, Korea)}
    {\textit{``Deep reinforcement learning for optimal mechanism in active Brownian particles''}}
    {Jul. 25, 2022}

\cvdatedentry{APCTP Workshop for Physics and Machine Learning (Jeju, Korea)}
    {\textit{``Exploring optimal mechanisms in active Brownian particles via deep reinforcement learning''}}
    {Nov. 26, 2021}

\cvdatedentry{Seoul National University Statistical Physics Seminar (Online)}
    {\textit{``Methods of estimating entropy production''}}
    {Feb. 1, 2021}

\pagebreak[4]
\cvdatedentry{Korean Physical Society Fall Meeting (Online)}
    {\textit{``Deep reinforcement learning for feedback-controlled flashing ratchets''}}
    {Nov. 6, 2020}

\cvdatedentry{NetSci2020 (Online)}
    {\textit{``Discovering wiring patterns of neural networks via backboning''}}
    {Sep. 22, 2020}

\cvdatedentry{Korean Physical Society Spring Meeting (Online)}
    {\textit{``Neural estimator for entropy production''}}
    {Jul. 13, 2020}

\cvdatedentry{Quantifying Success Satellite at NetSci2019 (Burlington, Vermont, USA)}
    {\textit{``Quantifying Individual Reputation in Large-scale Historical Documents''}}
    {May 27, 2019}

% === TECHNICAL SKILLS ===
\section{Technical Skills}

\begin{itemize}
    \item \textbf{Daily driver}: Currently running multiple Claude Code sessions (Max plan, \$200/mo) via tmux as my
primary research workflow.
    \item \textbf{ML/Interp}: PyTorch, JAX, TransformerLens, \& nnsight.
    \item \textbf{Infrastructure}: Distributed training \& inference (Slurm system with hundreds of H100s),
large-scale data pipelines
\end{itemize}

% === REVIEWING ===
\section{Academic Service}
\textbf{Program Committee}: AAAI 2027, CIKM 2026. \\
\textbf{Reviewer}: TMLR, NeurIPS 2026, ICML 2026 Workshop on Mechanistic Interpretability,
ICML 2026 (Received Gold Reviewer Award), ICLR 2026, NeurIPS 2025, ICML 2025 Workshop World Models,
ICLR 2025 Workshop Re-Align, NeurIPS 2024 Workshop Behavioral ML, ICML 2024 Workshop LLMs and Cognition.

% === AWARDS ===
\section{Awards}
\textbf{Pre-doctoral Fellow of Physics}, KAIST \hfill 2021

% === PRESS ===
\section{In the Press}
\begin{itemize}[itemsep=2pt]
    \item \textbf{Transformer as a hippocampal memory consolidation model based on NMDAR-inspired nonlinearity}
    (\textit{NeurIPS} 2023):
    \href{https://www.ibs.re.kr/cop/bbs/BBSMSTR_000000000738/selectBoardArticle.do?nttId=24264\&pageIndex=1\&searchCnd=\&searchWrd=}{IBS Research News (English)},
    \href{https://www.ibs.re.kr/cop/bbs/BBSMSTR_000000000735/selectBoardArticle.do?nttId=24270\&pageIndex=1\&searchCnd=\&searchWrd=}{IBS Research News (Korean)},
    \href{http://m.dongascience.com/news.php?idx=62694}{Donga Science}, and
    \href{https://www.eurekalert.org/news-releases/1029300}{EurekAlert!}.

    \item \textbf{Learning Entropy Production via Neural Networks}
    (\textit{Phys.\ Rev.\ Lett.} \textbf{125}, 140604, 2020):
    \href{https://webzine.kps.or.kr/?p=5_view\&idx=16495/}{Physics and High Technology}.
\end{itemize}

% === REFERENCES ===
\section{References}
\begin{multicols}{2}
\begin{description}[leftmargin=0pt,topsep=0pt,itemsep=5pt,parsep=0pt]
    \item[\href{https://www.mpi-sp.org/cha}{Meeyoung Cha}] \hfill \\
    Scientific Director \\
    Data Science for Humanity Group, MPI-SP \\
    \href{mailto:mia.cha@mpi-sp.org}{\faEnvelope \space mia.cha@mpi-sp.org}

    \item[\href{https://yongjoobaek.wordpress.com}{Yongjoo Baek}] \hfill \\
    Professor \\
    Department of Physics \& Astronomy, SNU \\
    \href{mailto:y.baek@snu.ac.kr}{\faEnvelope \space y.baek@snu.ac.kr}

    \item[\href{https://scholar.google.co.kr/citations?user=h1QXLx0AAAAJ\&hl=en}{Junghyo Jo}] \hfill \\
    Professor \\
    Department of Physics Education, SNU \\
    \href{mailto:jojunghyo@snu.ac.kr}{\faEnvelope \space jojunghyo@snu.ac.kr}
\end{description}
\columnbreak
\begin{description}[leftmargin=0pt,topsep=0pt,itemsep=5pt,parsep=0pt]
    \item[\href{https://stat.kaist.ac.kr:2583/~hjeong/}{Hawoong Jeong}] \hfill \\
    Professor \\
    Department of Physics, KAIST \\
    \href{mailto:hjeong@kaist.edu}{\faEnvelope \space hjeong@kaist.edu}

    \item[\href{https://centers.ibs.re.kr/html/glia_en/people/people_0201.html}{C.\ Justin Lee}] \hfill \\
    Director \\
    Center for Cognition and Sociality, IBS \\
    \href{mailto:cjl@ibs.re.kr}{\faEnvelope \space cjl@ibs.re.kr}
\end{description}
\end{multicols}

\end{document}

